\documentclass[11pt,a4paper]{article}
\usepackage[T1]{fontenc}
\usepackage{lmodern,microtype,amsmath,amssymb,amsthm,mathtools}
\usepackage[margin=24mm,headheight=15pt]{geometry}
\usepackage{booktabs,array,tabularx,xcolor,fancyhdr,titlesec}
\usepackage[colorlinks=true,linkcolor=teal!65!black,citecolor=teal!65!black,urlcolor=teal!65!black]{hyperref}
\definecolor{ink}{HTML}{15364A}
\definecolor{accent}{HTML}{007C78}
\titleformat{\section}{\Large\sffamily\bfseries\color{ink}}{\thesection}{.55em}{}
\titleformat{\subsection}{\large\sffamily\bfseries\color{accent}}{\thesubsection}{.55em}{}
\newtheorem{theorem}{Theorem}[section]
\newtheorem{lemma}[theorem]{Lemma}
\newtheorem{proposition}[theorem]{Proposition}
\newtheorem{corollary}[theorem]{Corollary}
\theoremstyle{remark}\newtheorem{remark}[theorem]{Remark}
\newtheorem{example}[theorem]{Example}
\DeclareMathOperator{\conv}{conv}
\DeclareMathOperator{\vol}{vol}
\DeclareMathOperator{\inter}{int}
\DeclareMathOperator{\relint}{relint}
\newcommand{\R}{\mathbb R}\newcommand{\Z}{\mathbb Z}
\pagestyle{fancy}\fancyhf{}
\fancyhead[L]{\small\sffamily Lattice coverings: remaining cases}
\fancyhead[R]{\small\sffamily Proof manuscript}
\fancyfoot[L]{\footnotesize\sffamily Prepared for Keio | 20 November 2026}
\fancyfoot[R]{\small\thepage}
\setlength{\parindent}{1.1em}\setlength{\parskip}{3pt}
\setlength{\emergencystretch}{2em}\raggedbottom
\hypersetup{pdftitle={Further Closures in Lattice Covering: Ambient Lattices, Interior Points, and Cartesian Products},pdfauthor={Research manuscript prepared for Keio}}
\begin{document}
\thispagestyle{empty}
\noindent{\sffamily\small\bfseries\color{accent} GEOMETRY OF NUMBERS | FOLLOW-UP TO VMRC [200]}\par\bigskip
\noindent{\sffamily\bfseries\fontsize{25}{29}\selectfont\color{ink}\raggedright
Further Closures in\\
Lattice Covering\par}
\medskip
\noindent{\large\sffamily Ambient lattices, interior points,\\
weighted minima, and Cartesian products}\par
\bigskip
\noindent Research manuscript prepared for Keio\hfill 20 November 2026
\bigskip
\begin{abstract}
We address four parts left open by the preceding work on the
Codenotti--Santos--Schymura covering problems. First, for every lattice
simplex containing an interior lattice origin, we give a necessary and
sufficient congruence criterion for covering radius $d/2$ in its actual
ambient lattice. This removes the vertex-generated-lattice restriction
from the earlier exact-covering-system criterion. Second, the numerical
interior-point conjecture is proved for every positive-integer-weight
standard terminal simplex, with its equality cases, and a three-nonprimitive-
vertex certificate treats additional polytopes. Third, for every
$n\ge2$ and real weights $0<a\le b\le\omega_j$ we determine exactly
$\mu_2(S_n(a,b,b,\omega_3,\ldots,\omega_n))$.
Fourth, Cartesian products of arbitrary convex bodies yield further
cases of the covering-product conjecture, including every product of
two or more factors of dimensions at most three. In particular, a
three-dimensional convex body times an interval has covering product
at least $1/3>5/16$. Arbitrary products of planar bodies have sharp
covering product at least $3/4$. Complete proofs are given, with existing
inputs and their elementary consequences distinguished from the new
scope of this note. The general numerical interior-point conjecture,
the general covering-product conjecture, the nonsimplicial equality
classification, and arbitrary weighted intermediate minima remain open
in this work.
\end{abstract}
\vfill
\noindent\begin{minipage}{\linewidth}
\small\color{ink}
\textbf{Scope.} This note records additional theorems beyond the preceding
23-page and 27-page transfer notes. It is not a replacement for those
notes and does not repeat their counterexamples. Mathematical correctness
is addressed by the proofs below; historical priority and external
peer review are separate questions.
\end{minipage}
\clearpage

\section{Targets, notation, and what changes}

For a full-dimensional compact convex body $K\subset\R^d$ and a
full-rank lattice $\Lambda$, its covering radius is
\[
 \mu(K,\Lambda)=\inf\{t\ge0:tK+\Lambda=\R^d\}.
\]
Its $i$th covering minimum is the least dilation whose lattice translates
meet every affine subspace of codimension $i$. We use the standard
quotient characterization \cite{CSS,GMS}:
\[
 \mu_i(K,\Lambda)=\sup_A\mu(AK,A\Lambda),
\]
where $A$ has rank $i$ and lattice-rational kernel. The minima increase
with $i$, and $\mu_d=\mu$. All lattices on sections and quotients are
the actual induced lattices. A non-hollow lattice polytope has at least
one interior lattice point.

For positive weights put
\[
 S_n(\omega)=\conv\{-\omega_0\mathbf1_n,
                    \omega_1e_1,\ldots,\omega_ne_n\},
 \qquad T_n=S_n(1,\ldots,1).
\]
The numerical part of CSS Conjecture D is
\begin{equation}\label{intro:D}
 \mu(P,\Lambda)\le\frac{d-1}{2}+\frac1{h+1},
 \qquad h=|\operatorname{int}P\cap\Lambda|\ge1.
\end{equation}
Conjecture F is
\begin{equation}\label{intro:F}
 \mathcal F_d(K,\Lambda):=
 \frac{\vol_dK}{\det\Lambda}\prod_{i=1}^{d}\mu_i(K,\Lambda)
 \ge\frac{d+1}{2^d}.
\end{equation}
For the ordered weights in the original Conjecture 5.3, the predicted
second minimum is
\[
 \frac{(\omega_0\omega_1)^{-1}+
       (\omega_0\omega_2)^{-1}+(\omega_1\omega_2)^{-1}}
      {\omega_0^{-1}+\omega_1^{-1}+\omega_2^{-1}}.
\]
That prediction is false in general by the preceding notes; here we
prove it on a surviving branch.

\begin{center}\small
\begin{tabularx}{\linewidth}{@{}>{\raggedright\arraybackslash}p{.16\linewidth}>{\raggedright\arraybackslash}X>{\raggedright\arraybackslash}X@{}}
\toprule
Target & Additional result in this note & Limit still present\\
\midrule
Equality in A & Complete congruence test for every simplex in any ambient lattice & General nonsimplicial polytopes\\
D, numerical & All integral-weight terminal simplices; structural certificates & Arbitrary lattice polytopes\\
Weighted minima & Exact $\mu_2$ when ordered weights begin $(a,b,b)$ & Arbitrary weights and higher indices\\
F & Cartesian classes with unrestricted convex factors; sharp planar-product constant & Arbitrary convex bodies\\
\bottomrule
\end{tabularx}
\end{center}

The weighted-simplex covering-radius formula and Schnell's planar
product inequality are published inputs. The universal constants used
for products are direct consequences of known descent arguments.
We include their derivations and make no assertion of historical
priority for these elementary combinations. The original numerical
non-hollow bound from \cite{VMRC200} is recalled with its proof below,
solely as an input to the additional results.

\section{Covering tools used in the proofs}

\begin{lemma}[Weighted fractional-part descent]\label{base:flow}
Let $q_j$ be positive integers, $c_j>0$, and $t_j\in\R$. Put
$C=\sum_jq_jc_j$ and $F(s)=\sum_jc_j\{t_j+sq_j\}$. Then
\[
 \min_s F(s)\le
 \frac12\sum_jc_j-\frac{\sum_jq_jc_j^2}{2C}.
\]
\end{lemma}
\begin{proof}
Between its finitely many events on a period, $F$ has slope $C$;
its mean is $\frac12\sum_jc_j$. At an event $e$ let $J_e$ be
the total downward jump and let $F(e)$ be the post-jump value.
Integrating the continuous derivative of $F^2$, and adding its
jumps, gives
\[
 \frac1C\sum_eJ_eF(e)
 =\frac12\sum_jc_j-\frac1{2C}\sum_eJ_e^2.
\]
Here $\sum_eJ_e=C$ and
$\sum_eJ_e^2\ge\sum_jq_jc_j^2$, since simultaneous positive
jumps only increase their squared total. The minimum is at most
the weighted average on the left.
\end{proof}

\begin{proposition}[Primitive-direction simplex bound]\label{base:simplex}
Let $S=\conv(v_0,\ldots,v_r)$ contain zero relatively internally,
write $v_j=\lambda_j u_j$ with $u_j$ primitive in its ambient lattice
and $\lambda_j>0$, and let $\alpha_j$ be the barycentric coordinates
of zero. Then
\[
 \mu(S)\le\frac12\sum_{j=0}^{r}\frac{1-\alpha_j}{\lambda_j}.
\]
\end{proposition}
\begin{proof}
The integer directions have a positive integer relation $q$.
Put $c_j=1/\lambda_j$ and $C=\sum_jq_jc_j$; then
$\alpha_j=q_jc_j/C$. Given $x=\sum_j t_j u_j$, subtract the lattice
vector $\sum_j\lfloor t_j+sq_j\rfloor u_j$. The remainder is
$\sum_jc_j\{t_j+sq_j\}v_j$, a nonnegative combination of the vertices.
Lemma~\ref{base:flow} bounds its coefficient sum by
$\frac12\sum_j(1-\alpha_j)c_j$. This proves the covering bound.
Only membership of the directions in the lattice is required, not
generation of that lattice.
\end{proof}

\begin{lemma}[Section and quotient]\label{base:section}
If $0\in\operatorname{int}K$ and $L$ is a lattice-rational linear
subspace with quotient map $\pi$, then
\[
 \mu(K,\Lambda)\le
 \mu(K\cap L,\Lambda\cap L)+\mu(\pi K,\pi\Lambda).
\]
\end{lemma}
\begin{proof}
Choose $a,b$ strictly above the two radii. For any $x$, use the
quotient covering to find $z\in\Lambda$ and $y\in bK$ with
$\pi(x-z)=\pi y$. Cover $x-z-y\in L$ by a translate of
$a(K\cap L)$ with a vector $w\in\Lambda\cap L$.
Then $x-z-w\in bK+aK=(a+b)K$. Let $a,b$ decrease to the radii.
\end{proof}

\begin{proposition}[Previously established non-hollow bound]\label{base:nonhollow}
Every non-hollow lattice $d$-polytope satisfies $\mu(P,\Lambda)\le d/2$.
\end{proposition}
\begin{proof}
Translate an interior lattice point to zero and choose a minimal set
of vertices whose convex hull contains zero. Its hull is a simplex
$S$ of dimension $r\ge1$ with zero in its relative interior.
Every primitive length of its lattice vertices is at least one,
so Proposition~\ref{base:simplex} gives $\mu(S)\le r/2$.
If $r=d$, inclusion proves the result. Otherwise its rational span
$L$ has section radius at most $r/2$, and the quotient $\pi P$ is
a non-hollow lattice polytope of dimension $d-r$. Induction and
Lemma~\ref{base:section} give $r/2+(d-r)/2$.
\end{proof}

\begin{proposition}[Established weighted terminal radius]\label{base:weighted}
For positive real weights,
\[
 \mu(S_d(\omega),\Z^d)=
 B(c):=\frac{\sum_{j<k}c_jc_k}{\sum_jc_j},
 \qquad c_j=1/\omega_j.
\]
\end{proposition}
\begin{proof}
This is the formula of \cite[Theorem 1.4]{CSS}. For completeness,
the directions $-\mathbf1_d,e_1,\ldots,e_d$ generate $\Z^d$ and
have relation $(1,\ldots,1)$. The same coefficient minimization
as in Lemma~\ref{eq:deep}, with weights $c_j$, gives the exact
distance $\min_s\sum_jc_j\{t_j+s\}$. Its upper bound from
Lemma~\ref{base:flow} is $B(c)$. For the reverse inequality,
arrange its $d+1$ events in a cyclic order with a gap
$c_j/\sum_kc_k$ preceding the jump of weight $c_j$. These gaps
sum to one. Every rise is cancelled by the following jump, so
all post-jump values coincide. The jump identity gives their
common value $B(c)$, proving equality.
\end{proof}


\section{A complete extremality test for a simplex in its ambient lattice}
\label{eq:section}

The vertex-generated hypothesis in the earlier exact-covering-system
criterion can be removed.  The replacement is a finite system of
pairwise congruence exclusions.  It is a recognition theorem for all
lattice simplices; it is not a classification of all extremal polytopes.

Let $S=\operatorname{conv}(v_0,\ldots,v_d)$ be a lattice simplex with
$0\in\operatorname{int}S$. Put
\[
 \Gamma=\sum_{j=0}^{d}\mathbb Z v_j\subseteq\Lambda,
 \qquad \sum_{j=0}^{d}q_jv_j=0,
 \qquad \gcd(q_0,\ldots,q_d)=1,
 \qquad Q=\sum_{j=0}^{d}q_j,
\]
where every $q_j$ is positive.  If each $q_j$ divides $Q$, set
$m_j=Q/q_j$. Define the facet forms by
\begin{equation}\label{eq:facetforms}
 \ell_j(v_i)=1-m_j\delta_{ij}.
\end{equation}
The positive relation makes these well-defined linear forms, and
$S=\{x:\ell_j(x)\le1\text{ for every }j\}$.

\begin{lemma}[Deep points of the coarse lattice]\label{eq:deep}
For $x=\sum_jt_jv_j$ one has
\begin{equation}\label{eq:distance}
 \rho_{S,\Gamma}(x)=\min_{s\in[0,1)}
                       \sum_{j=0}^{d}\{t_j+s q_j\}.
\end{equation}
The maximum of this function is at most $d/2$. It equals $d/2$
at $x$ if and only if the coordinate crossing events are simple and
form an equally spaced grid of $Q$ points on the circle.
Consequently a point at distance $d/2$ exists if and only if every $m_j$ is
an integer and there are residues $0\le a_j<m_j$ whose classes
$a_j+m_j\mathbb Z$ partition $\mathbb Z$. Its class modulo $\Gamma$
then has a representative
\begin{equation}\label{eq:deep-representative}
 x(a)=-\sum_j\frac{a_j}{m_j}v_j,
 \qquad \ell_j(x(a))=a_j-\frac d2.
\end{equation}
In particular, every point at distance $d/2$ satisfies
$\ell_j(x)\equiv-d/2\pmod{\mathbb Z}$ for every $j$.
\end{lemma}
\begin{proof}
Every translate by $\Gamma$ is obtained by replacing $t$ by an
integer translate, and every other coefficient representation has
the form $t+s q$. Minimizing the nonnegative coefficient sum gives
\eqref{eq:distance}. Write $F(s)=\sum_j\{t_j+s q_j\}$.
It has slope $Q$, mean $(d+1)/2$, and downward jumps $J_e$ at
the distinct events $e$. The jump identity, obtained by integrating
$d(F^2)/ds$ on the event intervals, is
\[
 \frac1Q\sum_e J_eF(e)
 =\frac{d+1}{2}-\frac1{2Q}\sum_eJ_e^2.
\]
Since $\sum_eJ_e=Q$ and every $J_e$ is a positive integer, the
right side is at most $d/2$. Equality of the minimum with $d/2$
forces every jump to be one and every post-jump value to be $d/2$.
The slope then forces gaps $1/Q$. The converse follows from the
same identity. A coordinate has crossings at spacing $1/q_j$,
so its grid positions are one residue class modulo $m_j=Q/q_j$.
They partition the grid exactly when they partition $\mathbb Z$.

For these classes, take $t_j=-a_j/m_j$. The resulting events
are exactly that grid. Exactly one residue $a_j$ is zero, so
$F(0)=d-\sum_j a_j/m_j=d/2$. Thus
$\sum_j a_j/m_j=d/2$, and applying \eqref{eq:facetforms} to
$x(a)$ gives \eqref{eq:deep-representative}. Conversely an arbitrary
point at distance $d/2$ has such coefficients after a common rotation of the
event parameter and changes by integer coefficients; the former
adds a multiple of $q$ and does not change the represented point.
The facet forms are integral on $\Gamma$, giving the last assertion.
\end{proof}

Choose any lattice basis $b_1,\ldots,b_d$ of $\Lambda$. When the
forms $\ell_j$ are integral on $\Lambda$, define
\begin{equation}\label{eq:Dij}
 D_{ij}=\gcd\bigl(m_i,m_j,
       (\ell_i-\ell_j)(b_1),\ldots,(\ell_i-\ell_j)(b_d)\bigr)
       \qquad(0\le i<j\le d).
\end{equation}
The gcd is positive and independent of the chosen lattice basis.

\begin{theorem}[Ambient-lattice congruence criterion]\label{eq:main}
The equality $\mu(S,\Lambda)=d/2$ holds if and only if all three
conditions below hold:
\begin{enumerate}
\item every $q_j$ divides $Q$;
\item every $\ell_j$ is integral on $\Lambda$;
\item there are integers $a_0,\ldots,a_d$, equivalently residues
      $0\le a_j<m_j$, such that
\begin{equation}\label{eq:avoid}
       a_i\not\equiv a_j\pmod{D_{ij}}
       \quad\text{for every }i<j.
\end{equation}
\end{enumerate}
If these conditions hold, \eqref{eq:deep-representative} gives a
deepest point for the full ambient lattice. In particular, a pair
with $D_{ij}=1$ rules out extremality.
\end{theorem}
\begin{proof}
Suppose first that $\mu(S,\Lambda)=d/2$. The distance is continuous
and lattice-periodic, so it has a maximizing point $x$. Since
$\Gamma\subseteq\Lambda$, for every $h\in\Lambda$ one has
\[
 d/2\le\rho_{S,\Gamma}(x-h)\le d/2.
\]
Lemma~\ref{eq:deep} gives condition (1). It also gives
$\ell_j(x),\ell_j(x-h)\equiv-d/2\pmod{\mathbb Z}$, hence
$\ell_j(h)\in\mathbb Z$, proving (2).
Replace $x$ by a representative $x(a)$ modulo $\Gamma$ as in
the lemma; this does not change its property of being deepest
for $\Lambda$.

We next describe the events for $x(a)-h$. For a coefficient
representation $y=\sum_jt_jv_j$, let $A=\sum_jt_j$. Then
$\ell_j(y)=A-m_jt_j$. In particular, if
$\ell_j(y)=r_j-d/2$ with all $r_j$ integral, the crossings of
$\{t_j+s q_j\}$ satisfy
\[
 Qs+A+d/2\equiv r_j\pmod{m_j}.
\]
Thus its distance equals $d/2$ if and only if the residue classes
$r_j+m_j\mathbb Z$ form an exact cover. Applied to $y=x(a)-h$,
these are $r_j=a_j-\ell_j(h)$.

Two residue classes with moduli $m_i,m_j$ intersect exactly when
their residues agree modulo $g_{ij}=\gcd(m_i,m_j)$. Their density
sum over all $j$ is $\sum_j1/m_j=1$. Since all the moduli divide
$Q$, pairwise disjointness is therefore equivalent to exact
coverage, by counting the residues in one period of length $Q$.
Consequently the translated systems are exact covers for every
$h\in\Lambda$ precisely when
\[
 a_i-a_j\notin
 \{(\ell_i-\ell_j)(h)\bmod g_{ij}:h\in\Lambda\}
 \quad\text{for every }i<j.
\]
The subgroup on the right is the set of multiples of $D_{ij}$
modulo $g_{ij}$, by the gcd definition and Bezout's identity.
This proves (3).

Conversely assume (1)--(3). Since $D_{ij}\mid g_{ij}$,
\eqref{eq:avoid} first implies that $a_j+m_j\mathbb Z$ is an exact
cover. Choose $x=x(a)$. The same subgroup argument shows that
$a_j-\ell_j(h)+m_j\mathbb Z$ remains an exact cover for every
$h\in\Lambda$. By the event calculation,
$\rho_{S,\Gamma}(x-h)=d/2$ for every such $h$. Finally
\[
 \rho_{S,\Lambda}(x)
 =\min_{h\in\Lambda/\Gamma}\rho_{S,\Gamma}(x-h)=d/2.
\]
The universal upper bound proves the desired equality.
\end{proof}

\begin{remark}[Relation to the earlier criterion]\label{eq:relation}
If $\Lambda=\Gamma$, the numbers in \eqref{eq:Dij} coming from
vertex generators are $0$, $-m_i$, or $m_j$. Hence
$D_{ij}=\gcd(m_i,m_j)$, and the theorem reduces to the exact
covering criterion of \cite{VMRC200}. For a strict superlattice
the smaller gcd values encode precisely the additional lattice
translates that may destroy a former deepest point. The theorem
is finite: after integral linear algebra one tests at most
$\prod_jm_j$ residue tuples. It makes no efficiency or polynomial
time claim.
\end{remark}

\begin{corollary}[Prime relation-sum rigidity]\label{eq:prime}
Suppose $Q=p$ is prime. Then $\mu(S,\Lambda)=d/2$ if and only if
all $q_j=1$ and $\Lambda=\Gamma$. In that case $d=p-1$, and
$S$ is a standard terminal simplex in its vertex-generated lattice.
\end{corollary}
\begin{proof}
Condition (1) forces every $q_j=1$, since $0<q_j<Q$. The vectors
$v_1,\ldots,v_d$ then form a basis of $\Gamma$ and
$v_0=-\sum_{j=1}^d v_j$. Use this basis to write $S=T_d$ and
$\Gamma=\mathbb Z^d$. We have $m_j=p$, and
\[
 \ell_0(x)=\sum_{k=1}^{d}x_k,
 \qquad \ell_j(x)=\ell_0(x)-p x_j.
\]
If $\Lambda$ strictly contains $\mathbb Z^d$, choose an $h\in\Lambda$
with a nonintegral coordinate $h_j$. Integrality of the facet
forms makes $p h_j$ an integer not divisible by $p$. Therefore
$D_{0j}=1$, contrary to condition (3). Conversely, in the standard
lattice the residues $0,1,\ldots,p-1$ give an exact cover.
\end{proof}

\begin{example}[A genuine superlattice that preserves extremality]
Let $S=T_3$ and
\[
 \Lambda=\mathbb Z^3+
          \mathbb Z(1/2,0,1/2).
\]
Here $m_0=m_1=m_2=m_3=4$ and the facet values of the added
generator are $(1,-1,1,-1)$. Thus $D_{02}=D_{13}=4$, and all
other $D_{ij}$ equal $2$. The residues $(a_0,a_1,a_2,a_3)=(0,1,2,3)$
satisfy every exclusion, proving $\mu(T_3,\Lambda)=3/2$.
The deepest representative is $(-1/4,-1/2,-3/4)$.
This is an illustration of the ambient-lattice test, not a new
counterexample to the known low-dimensional classification.
\end{example}

\begin{remark}[The scope of this closure]
Theorem~\ref{eq:main} decides the equality problem for every lattice
simplex with an interior lattice origin, including all ambient
superlattices. It does not assert that every extremal nonsimplicial
polytope decomposes into simplices of this type. The full corrected
polytope equality classification remains outside the theorem.
\end{remark}


\section{The interior-point inequality for every integral weighted terminal simplex}
\label{r4d:integral-terminal}

The exact radius of the weighted standard terminal simplex is an input
from Codenotti--Santos--Schymura, not a new radius formula.  The result
below compares that radius with the actual number of interior lattice
points and settles the numerical interior-point inequality throughout
the integral-weight family.  In particular, it includes simplices whose
interior lattice points are not collinear.

\begin{theorem}[All integral weights]
\label{r4d:weighted-D}
Let $d\ge1$, let $\omega_0,\ldots,\omega_d$ be positive integers, and put
\[
 S_d(\omega)=\operatorname{conv}
 \{-\omega_0\mathbf1_d,\omega_1e_1,\ldots,\omega_de_d\},
 \qquad h=|\operatorname{int}S_d(\omega)\cap\mathbb Z^d|.
\]
Then
\[
 \mu(S_d(\omega),\mathbb Z^d)
 \le \frac{d-1}{2}+\frac1{h+1}.
\]
For $d\ge2$, equality holds if and only if all $\omega_j=1$.
For $d=1$, equality holds for every pair of positive integer weights.
\end{theorem}
\begin{proof}
Write $c_j=1/\omega_j$, $C=\sum_jc_j$, and
\[
 B(c)=\frac{\sum_{j<k}c_jc_k}{C}
 =\frac C2-\frac{\sum_jc_j^2}{2C}.
\]
The known exact weighted-terminal formula gives $\mu(S_d(\omega))=B(c)$.
The function $B$ is increasing in every coordinate: with
$A=\sum_{j\ne i}c_j$ and $E=\sum_{j<k;\,j,k\ne i}c_jc_k$,
\[
 \frac{\partial B}{\partial c_i}
 =\frac{A^2-E}{(c_i+A)^2}>0
\]
when $d\ge1$. Permutations of the $d+1$ primitive directions
$-\mathbf1_d,e_1,\ldots,e_d$ are induced by unimodular transformations,
so the positions of the weights may be permuted.

For $d=1$, $S_1(\omega)=[-\omega_0,\omega_1]$, its number of interior
lattice points is $h=\omega_0+\omega_1-1$, and its covering radius is
$1/(h+1)$.  Henceforth assume $d\ge2$, and let $r$ be the number of
weights larger than one.

If $r=0$, this is the standard terminal simplex, with $h=1$ and
$\mu=d/2$.

If $r=1$, move the enlarged weight $a\ge2$ to the first positive
coordinate. Every interior lattice point has all other coordinates
zero and first coordinate $0,1,\ldots,a-1$. Thus $h=a$, while
\[
 B(c)=\frac{d-1}{2}+\frac{d+1}{2(da+1)}.
\]
The exact difference from the desired bound is
\[
 \frac1{a+1}-\frac{d+1}{2(da+1)}
 =\frac{(d-1)(a-1)}{2(a+1)(da+1)}>0.
\]

If $r=2$, move the enlarged weights to the first two positive
coordinates and call them $a,b\ge2$. All remaining coordinates of an
interior lattice point vanish. The first two coordinates are exactly
the nonnegative integers $(x,y)$ such that $x/a+y/b<1$. Consequently
\[
 h=\frac{ab+a+b-\gcd(a,b)}2.
\]
This count also follows by applying Pick's formula to the triangle
with vertices $(0,0),(a,0),(0,b)$ and removing its slanted edge.
Put $m=d-1$. Direct simplification gives
\[
 B(c)=\frac m2+
 \frac{m(a+b-ab)+2}{2(mab+a+b)}.
\]
If $a=b=2$, then $h=3$ and the last fraction is $1/(4d)<1/4$.
Otherwise assume $a\le b$. If $(a,b)=(2,3)$, the numerator is $2-m$:
it is nonpositive for $d\ge3$, while for $d=2$ the fraction is
$1/22<1/6=1/(h+1)$ because $h=5$.
Every other pair has $a+b-ab\le-2$, making the numerator nonpositive.
In all these cases the desired inequality is strict.

Finally, suppose $r\ge3$. By monotonicity, replacing three enlarged
weights by $2$ and all other weights by $1$ can only increase $B$.
For the resulting reciprocal weights, three entries are $1/2$ and
the remaining $d-2$ entries are $1$. Therefore
\[
 B(c)\le
 \frac{d-1}{2}+\frac{2-d}{4d-2}
 \le\frac{d-1}{2}
 <\frac{d-1}{2}+\frac1{h+1}.
\]
This proves the inequality and all equality assertions.
\end{proof}

\begin{corollary}[A contained-simplex certificate independent of the interior count]
\label{r4d:witness}
Let $P$ be a lattice $d$-polytope with $0\in\operatorname{int}P$.
Suppose a rational $r$-dimensional linear subspace $L$, $r\ge2$,
contains a lattice $r$-simplex $T\subseteq P\cap L$ with
$0\in\operatorname{relint}T$ and at least three vertices nonprimitive
in $\Lambda\cap L$.
Then
\[
 \mu(P,\Lambda)\le\frac{d-1}{2}.
\]
In particular, $P$ satisfies the numerical interior-point conjecture,
strictly, regardless of its number of interior lattice points.
\end{corollary}
\begin{proof}
Write the vertices as $v_j=\lambda_j u_j$, where $u_j$ is primitive,
$\lambda_j$ is a positive integer, and $\alpha_j$ is the barycentric
coordinate of zero at $v_j$. Choose three indices $J$ with
$\lambda_j\ge2$. The primitive-direction surface bound gives
\[
 \mu(T,\Lambda\cap L)
 \le\frac r2-\frac14\sum_{j\in J}(1-\alpha_j)
 =\frac r2-\frac34+\frac14\sum_{j\in J}\alpha_j
 \le\frac{r-1}{2}.
\]
If $r=d$, monotonicity under inclusion proves the claim.
Otherwise the projection of $P$ to $\mathbb R^d/L$ is a non-hollow
lattice polytope. Apply the section-and-quotient inequality and the
non-hollow bound to obtain
\[
 \mu(P,\Lambda)
 \le\frac{r-1}{2}+\frac{d-r}{2}=\frac{d-1}{2}.
\]
\end{proof}

More generally, the same argument gives
$\mu(T)\le r/2-(s-1)/4$ if $T$ has $s$ nonprimitive vertices.
The three-vertex certificate does not require a terminal-simplex
relation, a vertex-generated ambient lattice, or any upper bound on
the total number of interior lattice points.

\paragraph{Limits of this closure.}
The theorem settles the numerical interior-point inequality for all
integral weighted standard terminal simplices in the standard ambient
lattice. It does not settle the inequality for every lattice simplex,
for arbitrary finer ambient lattices, or for arbitrary polytopes.
The difficult two-interior-point endpoint regime from the previous
manuscript remains unresolved. The exact radius input was already known;
no historical priority is asserted for this new comparison without a
separate literature audit.


% New research fragment, 20 November 2026.
% The proof uses only elementary planar lattice refinement, fractional-part
% rounding, section/projection subadditivity, and the established endpoint
% formula of CSS. Labels introduced here start with r4w:.

\section{A surviving exact branch of the weighted covering-minima problem}
\label{r4w:sec-one-small}

The preceding counterexamples to the weighted covering-minima conjecture
do not determine all of its surviving parameter branches. We now determine
the second minimum when the two weights after the smallest one coincide.
In particular, this settles the one-small-weight family in every ambient
dimension and for every real value of its parameter.

\begin{lemma}[Three first-event bounds]
\label{r4w:lem-first-events}
Let $p,q\ge1$ be integers and $0<b\le1$. For all real $t_0,t_1,t_2$,
\[
 \min_{s\in\mathbb R}
 \bigl(\{t_0+s\}+b\{t_1+ps\}+b\{t_2+qs\}\bigr)
 \le \frac{b(2+b)}{1+2b}.
\]
\end{lemma}
\begin{proof}
Rotate the parameter to make $t_0=0$ modulo one, and replace the other
phases by their representatives in $[0,1)$. Their first subsequent
crossing times are
\[
 \alpha=\frac{1-t_1}{p},\qquad
 \beta=\frac{1-t_2}{q}.
\]
Interchanging the two equally weighted coordinates, if needed, makes
$\alpha\le\beta$. Both times belong to $(0,1]$. Put
$C=1+b(p+q)$ and denote the displayed fractional-part sum by $F$.
Then
\begin{align*}
 F(0)&=L_0:=2b-bp\alpha-bq\beta,\\
 F(\alpha)&\le L_1:=b+(1+bq)\alpha-bq\beta,\\
 F(\beta)&\le L_2:=(1+bp)\beta-bp\alpha.
\end{align*}
The second estimate subtracts at least the first $p$-coordinate jump.
The third subtracts at least one jump of each of the two coordinates.
Any extra crossings only strengthen these inequalities, including
coincidences and a crossing of the first coordinate at time one.

The positive numbers $1/C,bp/C,bq/C$ sum to one. Their corresponding
weighted average of $L_0,L_1,L_2$ is
\[
 \frac{L_0+bpL_1+bqL_2}{C}
 =\frac{b(2+bp)}{1+b(p+q)}.
\]
Thus the minimum of $F$ is at most this last quantity. Finally,
\[
 \frac{b(2+b)}{1+2b}
 -\frac{b(2+bp)}{1+b(p+q)}
 =\frac{b^2\bigl((1-b)(p-1)+(2+b)(q-1)\bigr)}
 {(1+2b)(1+b(p+q))}\ge0.
\]
This proves the lemma without a restriction on later crossings.
\end{proof}

\begin{lemma}[Planar unimodular refinement inside a triangle]
\label{r4w:lem-fan}
Let $v,w$ be independent vectors of a planar lattice $\Lambda$.
Their cone admits a subdivision into unimodular cones whose primitive
ray generators all lie in $\operatorname{conv}(0,v,w)$.
\end{lemma}
\begin{proof}
First shorten the two boundary vectors to their primitive generators.
Suppose a current pair of primitive generators $u_1,u_2$ has lattice
determinant $D>1$. A nonzero lattice class in their half-open
fundamental parallelogram has a representative
$z=\alpha u_1+\beta u_2$ with $0\le\alpha,\beta<1$.
Primitivity of the two generators forces $0<\alpha,\beta<1$.
If $\alpha+\beta\le1$, retain $z$; otherwise retain $u_1+u_2-z$.
The retained vector belongs to $\operatorname{conv}(0,u_1,u_2)$
and lies strictly between the boundary rays. Shorten it to its
primitive generator and subdivide along that ray.
Both new positive lattice determinants are strictly smaller than $D$.
Induction terminates in determinant-one cones. Every new triangle
lies in its predecessor, so all generators remain in the original
triangle.
\end{proof}

\begin{theorem}[Planar stretching with at most one fixed column]
\label{r4w:thm-stretch}
Let $u_1,\ldots,u_N$ be vectors of a planar lattice $\Lambda$,
with $0\in\operatorname{int}\operatorname{conv}(u_1,\ldots,u_N)$.
Keep at most one of the columns fixed and multiply all the others
by $r\ge1$. For their resulting convex hull $P_r$,
\[
 \boxed{\displaystyle
 \mu(P_r,\Lambda)\le G(r):=\frac{2r+1}{r(r+2)}.}
\]
The columns need not be primitive, distinct, or generators of $\Lambda$.
\end{theorem}
\begin{proof}
The case $r=1$ is included below. Delete zero columns, which does not
affect origin-interiority or the convex hull. A minimally supported
positive dependence among the columns consists either of an opposite
pair or of a triangle containing zero in its interior.

First consider such a triangle. If none of its vertices is fixed,
its scaled hull has covering radius at most $1/r\le G(r)$ by the
planar non-hollow bound. If exactly one is fixed, write its hull as
\[
 T_r=\operatorname{conv}(v_0,rv_1,rv_2),
 \qquad 0\in\operatorname{int}\operatorname{conv}(v_0,v_1,v_2).
\]
The vectors $v_1,v_2$ are independent, and $-v_0$ lies strictly
inside their positive cone. Apply Lemma~\ref{r4w:lem-fan} to that cone.

If $-v_0$ lies strictly inside one of its unimodular cones, let
$w_1,w_2$ be the generators of that cone. Then
\[
 v_0=-pw_1-qw_2\quad\text{for integers }p,q\ge1,
\]
and
$T'_r=\operatorname{conv}(v_0,rw_1,rw_2)\subseteq T_r$.
The inclusion holds because $rw_j\in\operatorname{conv}(0,rv_1,rv_2)$
and $0\in T_r$. For any point written in the integer directions
$v_0,w_1,w_2$, fractional-part rounding along their relation
$(1,p,q)$ bounds its distance to $\Lambda$ by
\[
 \min_s\bigl(\{t_0+s\}+r^{-1}\{t_1+ps\}
                    +r^{-1}\{t_2+qs\}\bigr).
\]
Lemma~\ref{r4w:lem-first-events}, with $b=1/r$, bounds this by
$G(r)$, as required.

If $-v_0$ lies on a ray of the subdivision, write $v_0=-Aw$
with $w$ primitive and $A\ge1$ an integer. The section of $T_r$
along $\mathbb Rw$ contains $[-Aw,rw]$, of lattice length at
least $r+1$. The quotient of $T_r$ by this line has lattice
length at least $2r$: its two nonzero endpoint candidates are
$r$ times opposite-sign nonzero integers. Section/projection
subadditivity therefore gives
\[
 \mu(T_r,\Lambda)\le\frac1{r+1}+\frac1{2r}\le G(r),
\]
where
\[
 G(r)-\left(\frac1{r+1}+\frac1{2r}\right)
 =\frac{r-1}{2(r+1)(r+2)}\ge0.
\]

It remains to handle an opposite pair in the original columns.
Let $W$ be its line. If one selected column is fixed, the section
of $P_r$ along $W$ has lattice length at least $r+1$, and all
nonzero projected columns are scaled, so the quotient interval
has lattice length at least $2r$. If both selected columns are
scaled, the section length is at least $2r$. At most one nonzero
projected column is fixed; because the unscaled quotient has
opposite-sign nonzero integer endpoints, the scaled quotient has
length at least $r+1$. Thus both cases give the same preceding
bound. Each selected triangle or section lies in $P_r$, so
monotonicity completes the proof.
\end{proof}

\begin{theorem}[Exact second minimum with one smallest weight]
\label{r4w:thm-mu2}
For every integer $n\ge2$ and real $r\ge1$,
\[
 \boxed{\displaystyle
 \mu_2\bigl(S_n(1,r,\ldots,r),\mathbb Z^n\bigr)
 =\frac{2r+1}{r(r+2)}.}
\]
More generally, for $0<a\le b$ and arbitrary
$\omega_3,\ldots,\omega_n\ge b$,
\[
 \boxed{\displaystyle
 \mu_2\bigl(S_n(a,b,b,\omega_3,\ldots,\omega_n),\mathbb Z^n\bigr)
 =\frac{a+2b}{b(2a+b)}.}
\]
For $n=2$, the list of extra weights is empty.
\end{theorem}
\begin{proof}
For any rational rank-two quotient $A$, its unweighted columns
$A(-\mathbf1_n),Ae_1,\ldots,Ae_n$ belong to the quotient lattice
$A\mathbb Z^n$, span the plane, and sum to zero. Consequently
zero is interior to their convex hull. The image of
$S_n(1,r,\ldots,r)$ keeps just the first column fixed and scales
all remaining columns by $r$. Theorem~\ref{r4w:thm-stretch} gives
the uniform quotient-radius bound $G(r)$.
The standard rational-quotient characterization of covering minima
therefore implies $\mu_2\le G(r)$.

The coordinate quotient onto any two standard coordinates is
$S_2(1,r,r)$. The established weighted-simplex radius formula gives
\[
 \mu(S_2(1,r,r))
 =\frac{2/r+1/r^2}{1+2/r}=G(r),
\]
proving the matching lower bound.

For the more general weights, the simplex contains
$S_n(a,b,b,b,\ldots,b)$, so monotonicity and scaling supply the
same upper bound. Its first-two-coordinate quotient is still
$S_2(a,b,b)$, supplying the same lower bound.
\end{proof}

\begin{corollary}[The complete covering vector of one-small-weight tetrahedra]
\label{r4w:cor-vector}
For $0<a\le b$,
\[
 (\mu_1,\mu_2,\mu_3)(S_3(a,b,b,b))
 =\left(\frac1{a+b},\frac{a+2b}{b(2a+b)},
              \frac{3(a+b)}{b(3a+b)}\right).
\]
\end{corollary}
\begin{proof}
The middle component is Theorem~\ref{r4w:thm-mu2}; the first and
last are the established width and weighted covering-radius formulas.
\end{proof}

\begin{remark}[Exact scope]
This proves the coordinate-projection prediction for $i=2$ in the
entire one-small-weight family, and in the larger family with ordered
weights beginning $(a,b,b)$. It does not restore the original conjecture
for arbitrary weights or for higher intermediate indices. In fact the
same one-small-weight family has counterexamples to that prediction
at every intermediate index $i\ge3$, as established in the preceding
manuscript. Planar unimodular refinement inside the original triangle
is the step that makes the present argument specifically planar.
\end{remark}


% Round 4 fragment. Required bibliography keys: GMS, Schnell.
% GMS: Gonzalez Merino and Schymura, DCG 58 (2017), 663--685.
% An accessible author preprint is OWP 2016-08 (under Henze's former name).
% No unproved case of Conjecture F is used below.

\section{Covering products of arbitrary Cartesian factors}
\label{r4f:section}

Write
\[
 V(K,\Lambda)=\frac{\operatorname{vol}_d(K)}{\det\Lambda},\qquad
 \mathcal F_d(K,\Lambda)=V(K,\Lambda)\prod_{i=1}^d\mu_i(K,\Lambda),
 \qquad \theta(K,\Lambda)=V(K,\Lambda)\mu_d(K,\Lambda)^d.
\]
The elementary covering-density bound is $\theta(K,\Lambda)\ge1$.
The product constructions in this section use product lattices;
arbitrary simultaneous invertible linear images of the body and lattice
are allowed. Cartesian products, denoted by $\times$, are distinct
from the free sums used in the preceding transfer manuscript.

\begin{theorem}[Exact covering minima of Cartesian products]
\label{r4f:cartesian-spectrum}
Let $K_j\subset\mathbb R^{d_j}$ be arbitrary convex bodies, let
$\Lambda_j$ be full-rank lattices, and put
$P=\prod_{j=1}^m K_j$, $\Lambda=\bigoplus_j\Lambda_j$, and
$D=\sum_jd_j$. Then, for every $1\le i\le D$,
\begin{equation}
\label{r4f:exact-spectrum}
 \mu_i(P,\Lambda)
 =\max_{1\le j\le m}\mu_{\min\{i,d_j\}}(K_j,\Lambda_j).
\end{equation}
No symmetry, integrality, or interior-lattice-point hypothesis is imposed
on the factors.
\end{theorem}

\begin{proof}
Use the characterization of covering minima as suprema of covering
radii of rational linear quotients. For the upper bound, let $A$ be any
rank-$i$ quotient of the product lattice, and let $A_j$ be its restriction
to the $j$th factor. Put $U_j=\operatorname{im}A_j$ and $r_j=\dim U_j$.
The subgroup $A_j\Lambda_j$ is discrete, since it is a subgroup of the
lattice $A\Lambda$, and it is a full-rank lattice in $U_j$.
For a zero-dimensional $U_j$ the covering radius is understood to be zero.
For every
\[
 t>\max_j\mu_{\min\{i,d_j\}}(K_j,\Lambda_j),
\]
the quotient characterization and monotonicity give
$tA_jK_j+A_j\Lambda_j=U_j$ for each $j$. Adding these identities yields
\[
 tAP+A\Lambda
 =\sum_j(tA_jK_j+A_j\Lambda_j)=\sum_jU_j=\mathbb R^i.
\]
Take the supremum over $A$ and then let $t$ decrease.

For the reverse bound, fix $j$ and put $r=\min\{i,d_j\}$.
Choose any rank-$r$ quotient of the $j$th factor and complement it with
a rank-$(i-r)$ quotient of the other factors. Such a complement exists
because $i-r\le D-d_j$. The resulting quotient body and lattice are
Cartesian products. The covering radius of a Cartesian product is the
maximum of the radii of its factors, directly from the definition.
Thus $\mu_i(P,\Lambda)$ is at least the radius of the chosen quotient
of $K_j$. Taking its supremum proves the required lower bound.
\end{proof}

\begin{lemma}[A standard projection descent]
\label{r4f:descent}
For every $d$-dimensional convex body $K$ with $d\ge2$, there is a
rational quotient $Q$ of dimension $d-1$ such that
\[
 V(K,\Lambda)\ge\frac{V(Q,\Lambda_Q)}{d\mu_d(K,\Lambda)},
 \qquad
 \mu_i(Q,\Lambda_Q)\le\mu_i(K,\Lambda)\quad(1\le i<d).
\]
\end{lemma}

\begin{proof}
This is the descent in the proof of the general $1/d!$ bound in
\cite{GMS}; details are included to specify exactly the input being used.
The standard inequality $\lambda_1(K-K,\Lambda)\le\mu_d(K,\Lambda)$
provides a nonzero lattice vector
$z=\mu_d(K,\Lambda)(x-y)$ with $x,y\in K$.
For the rational line $L=\mathbb Rz$, the body therefore has a chord
parallel to $L$ of length at least
$\det(\Lambda\cap L)/\mu_d(K,\Lambda)$.
The elementary section--projection volume inequality states that the
volume of $K$ is at least $1/d$ times the length of any such chord
times the $(d-1)$-volume of $Q=K|L^\perp$.
One may see this by using concavity of the lengths of the parallel
fibres: the cone under the maximal fibre has integral equal to $1/d$
times its height times the volume of its projection.
Divide by
$\det\Lambda=\det(\Lambda\cap L)\det(\Lambda|L^\perp)$.
The covering-minimum comparison follows by composing rational quotients.
\end{proof}

\begin{lemma}[Suffix products and the planar starting constant]
\label{r4f:suffix}
For $1\le r\le d$,
\begin{equation}
\label{r4f:suffix-bound}
 V(K,\Lambda)\left(\prod_{i=r}^d\mu_i(K,\Lambda)\right)
                  \mu_d(K,\Lambda)^{r-1}\ge\frac{r!}{d!}.
\end{equation}
Moreover,
\begin{equation}
\label{r4f:beta}
 \mathcal F_d(K,\Lambda)\ge\beta_d,
 \qquad
 \beta_1=1,\quad\beta_d=\frac{3}{2d!}\quad(d\ge2).
\end{equation}
The constant in \eqref{r4f:beta} is a direct consequence of the
published projection descent and Schnell's planar theorem; it is not
presented as a new difficult universal inequality.
\end{lemma}

\begin{proof}
For $d=r$, \eqref{r4f:suffix-bound} is the covering-density bound.
For $d>r$, apply Lemma~\ref{r4f:descent}; cancel one factor of
$\mu_d(K,\Lambda)$, compare the remaining indexed factors with the
covering minima of $Q$, and use
$\mu_d(K,\Lambda)\ge\mu_{d-1}(Q,\Lambda_Q)$.
The resulting expression is at least $1/d$ times the same suffix
product in dimension $d-1$. Iteration gives $r!/d!$.
For \eqref{r4f:beta}, use Schnell's planar bound
$\mathcal F_2\ge3/4$ \cite{Schnell}, also recorded in \cite{GMS},
and iterate the same descent from dimension $d$ down to dimension two.
The accumulated factor is $2/d!$, proving the formula.
\end{proof}

\begin{theorem}[A product bound requiring no conjecture for the factors]
\label{r4f:unconditional-products}
Under the hypotheses of Theorem~\ref{r4f:cartesian-spectrum}, let
$p=\max_jd_j$ and $M=\max_j\mu_{d_j}(K_j,\Lambda_j)$.
For every $j_*$ with $d_{j_*}=p$,
\begin{equation}
\label{r4f:strong-product}
 \mathcal F_D(P,\Lambda)
 \ge\mathcal F_p(K_{j_*},\Lambda_{j_*})
       \prod_{j\ne j_*}\theta(K_j,\Lambda_j)
 \ge\beta_p.
\end{equation}
Consequently Conjecture F holds for every such product whenever
\begin{equation}
\label{r4f:threshold}
 \beta_p\ge\frac{D+1}{2^D}.
\end{equation}
\end{theorem}

\begin{proof}
Put $c_i=\mu_i(P,\Lambda)$. The exact formula gives $c_i=M$ for
$p\le i\le D$, while $c_i\ge\mu_i(K_{j_*},\Lambda_{j_*})$
for $1\le i\le p$. Since normalized Cartesian-product volumes
multiply and $D-p=\sum_{j\ne j_*}d_j$, it follows that
\begin{align*}
 \mathcal F_D(P,\Lambda)
 &\ge \mathcal F_p(K_{j_*},\Lambda_{j_*})
          \prod_{j\ne j_*} V(K_j,\Lambda_j)M^{d_j}\\
 &\ge \mathcal F_p(K_{j_*},\Lambda_{j_*})
          \prod_{j\ne j_*}\theta(K_j,\Lambda_j).
\end{align*}
Apply the density bound and Lemma~\ref{r4f:suffix}.
\end{proof}

For example, the least $D\ge p$ for which \eqref{r4f:threshold}
holds is
\[
\begin{array}{c|rrrrrrr}
p&2&3&4&5&6&7&8\\\hline
D&2&5&7&10&13&16&20.
\end{array}
\]
Thus any pair of arbitrary three-, four-, or five-dimensional convex
bodies of the same dimension has a Cartesian product satisfying F.
No assertion that F holds for either individual factor is used.

\begin{theorem}[Products with arbitrary rectangular boxes]
\label{r4f:boxes}
Let $p,q\ge1$ be integers, let $K\subset\mathbb R^p$ be arbitrary, and let
$B=I_1\times\cdots\times I_q\subset\mathbb R^q$ be an axis-parallel
box, where each interval has positive length and the lattice is the
product lattice. Then
\begin{equation}
\label{r4f:box-bound}
 \mathcal F_{p+q}(K\times B)\ge
 \frac{\min\{p,q+1\}!}{p!}.
\end{equation}
In particular, for every three-dimensional $K$ and every interval $I$,
\begin{equation}
\label{r4f:three-prism}
 \mathcal F_4(K\times I)\ge\frac13>\frac5{16}.
\end{equation}
\end{theorem}

\begin{proof}
Normalize the interval lattices to $\mathbb Z$, write
$s=\max_j |I_j|^{-1}$, and set $a_i=\mu_i(K)$.
Then $\operatorname{vol}_q(B)\ge s^{-q}$ and every covering minimum
of $B$ is $s$. The exact Cartesian formula yields
\begin{equation}
\label{r4f:box-objective}
 \mathcal F_{p+q}(K\times B)
 \ge V(K)s^{-q}
       \left(\prod_{i=1}^{p}\max\{a_i,s\}\right)
       \max\{a_p,s\}^{q}.
\end{equation}
If $q<p$, the right-hand side, minimized over $s>0$, equals
\[
 V(K)\left(\prod_{i=q+1}^{p}a_i\right)a_p^q.
\]
Indeed, for $s\le a_p$, its exponent of $s$ between consecutive
$a_i$ is $\#\{i:a_i<s\}-q$; a minimum occurs on
$[a_q,a_{q+1}]$. For $s\ge a_p$ it is $V(K)s^p$ and increases.
If $q\ge p$, the minimum is $V(K)a_p^p$, attained at $s=a_p$.
Now use \eqref{r4f:suffix-bound} with $r=q+1$ in the first case
and the density bound in the second case. This proves
\eqref{r4f:box-bound}; the displayed four-dimensional consequence
uses $p=3$, $q=1$.
\end{proof}

\begin{corollary}[All nontrivial Cartesian products of low-dimensional bodies]
\label{r4f:low-dimensional}
Let $m\ge2$, and let each $K_j$ be an arbitrary convex body of
dimension $d_j\in\{1,2,3\}$. Then $P=\prod_jK_j$, with the product
lattice, satisfies Conjecture F.
\end{corollary}

\begin{proof}
If $p\le2$, Theorem~\ref{r4f:unconditional-products} suffices.
If $p=3$ and $D\ge5$, it gives $\mathcal F_D(P)\ge1/4$,
while $(D+1)/2^D\le6/32<1/4$.
The only remaining case is $D=4$, necessarily a three-dimensional
factor times an interval, covered by \eqref{r4f:three-prism}.
\end{proof}

\begin{theorem}[Sharp bound for arbitrary products of planar bodies]
\label{r4f:planar-products}
For arbitrary planar convex bodies $K_1,\ldots,K_m$ and their
product lattice,
\[
 \mathcal F_{2m}(K_1\times\cdots\times K_m)\ge\frac34.
\]
The constant $3/4$ is attained for every $m\ge1$.
\end{theorem}

\begin{proof}
The inequality is \eqref{r4f:strong-product} with $p=2$.
For sharpness, use the hexagon
\[
 H=\operatorname{conv}\left\{\pm(2/3,-1/3),
               \pm(1/3,1/3),\pm(-1/3,2/3)\right\}.
\]
Its area is one. The inequalities defining it are
$|2x+y|\le1$, $|x+2y|\le1$, and $|x-y|\le1$.
The interior of $2H$ contains no nonzero integer vector: the three
strict inequalities at an integer pair force each absolute integer
expression to be at most one, which is possible only at the origin.
Thus the lattice translates of $H$ have disjoint interiors. Their
density is one, so they cover: otherwise the complement of their
locally finite closed union would contain an open set of positive area.
It follows that $\mu_2(H)=1$.
For a nonzero integer covector $(u,v)$ the width of $H$ is
\[
 \frac23\max\{|2u-v|,|u+v|,|-u+2v|\}.
\]
The integer maximum is at least two and equals two for $(u,v)=(1,0)$.
Indeed, writing $A=2u-v$, $B=u+v$, the condition
$|A|,|B|,|B-A|\le1$ and $u=(A+B)/3\in\mathbb Z$
forces $A=B=0$. Hence $\mu_1(H)=3/4$.
The exact product formula gives the spectrum
$(3/4,1,\ldots,1)$ for $H^m$. Its volume is one, proving sharpness.
\end{proof}

\paragraph{Scope and attribution.}
These results settle F on the stated Cartesian-product classes, which
are not restricted to the earlier free sums of weighted terminal
simplices. The auxiliary universal constants are consequences of the
published GMS descent and Schnell's planar theorem. The direct
Cartesian formula is elementary and is proved here rather than claimed
as a historically new principle. None of the results gives F
for an arbitrary indecomposable convex body, nor transfers the product
result back to an individual factor.


\section{What remains after these additional closures}

The ambient-lattice restriction has been removed completely for the
simplex equality problem: Theorem~\ref{eq:main} has no generation
hypothesis and is a genuine necessary and sufficient finite criterion.
Its congruence formulation does not classify every nonsimplicial
polytope with radius $d/2$.

For the numerical interior-point problem, Theorem~\ref{r4d:weighted-D}
handles every integral weight vector in every dimension and determines
all equality cases in that family. The nonprimitive-vertex certificate
also applies to polytopes containing a suitable simplex in a rational
section. Neither statement closes the earlier missing two-point regime.
If the line through two interior lattice points has normalized section
\[
 [-\beta u,(1+\delta)u],\qquad 0<\beta,\delta\le1,
\]
the previous method still leaves
$2\max(\beta,\delta)<1+\min(\beta,\delta)$ untreated in general.
In particular, the new weighted-terminal family theorem should not be
described as resolving that entire two-point gap.

Theorem~\ref{r4w:thm-mu2} is an exact formula, not just a lower bound
from a selected projection. It applies to all real parameters in its
specified branch and all ambient dimensions. The general weighted
intermediate-minimum problem remains, and the earlier counterexamples
at higher indices remain valid.

The Cartesian-product results impose no unproved version of F on the
individual factors. They therefore enlarge the proven classes beyond
the earlier free-sum constructions. They do not imply F for an arbitrary
factor by reversing the product construction. General convex bodies
outside the stated product classes remain beyond the proof.

\subsection*{Verification and dependence}
Each additional argument was read independently during preparation.
The checks specifically included simultaneous and endpoint events in
the planar flow, the ambient-lattice quantifiers in the congruence test,
the actual interior-point count in the two-enlarged-weight case, and
the indices and volume normalization in the product descent.
The displayed theorems are proved analytically; finite enumeration is
not used as a premise. This manuscript is not a formal proof-assistant
certificate and does not assert external peer-review acceptance or
historical priority.

\begin{thebibliography}{9}
\interlinepenalty=10000
\raggedright
\bibitem{CSS}
G.~Codenotti, F.~Santos and M.~Schymura,
\emph{The covering radius and a discrete surface area for non-hollow
simplices}, Discrete \& Computational Geometry \textbf{67} (2022),
65--111.
\href{https://doi.org/10.1007/s00454-021-00330-3}{doi:10.1007/s00454-021-00330-3};
\href{https://arxiv.org/abs/1903.02866v2}{arXiv:1903.02866v2}.

\bibitem{GMS}
B.~Gonz\'alez Merino and M.~Schymura,
\emph{On densities of lattice arrangements intersecting every
$i$-dimensional affine subspace}, Discrete \& Computational Geometry
\textbf{58} (2017), 663--685.
\href{https://doi.org/10.1007/s00454-017-9911-x}{doi:10.1007/s00454-017-9911-x}.
Accessible author preprint, under the author name M.~Henze:
\href{https://publications.mfo.de/bitstream/handle/mfo/1118/OWP2016_08.pdf?isAllowed=y\&sequence=1}
{Oberwolfach Preprint 2016-08}; see Theorems 4.1 and 4.3.

\bibitem{Schnell}
U.~Schnell, \emph{A Minkowski-type theorem for covering minima in
the plane}, Geometriae Dedicata \textbf{55} (1995), 247--255.
\href{https://doi.org/10.1007/BF01266316}{doi:10.1007/BF01266316}.

\bibitem{VMRC200}
\emph{Weighted fractional-part flows and sharp lattice covering bounds},
Verifiable Mathematical Reasoning Corpus, manuscript [200].
Repository revision
\texttt{ac1492a50784be919fe0b5e8f19f96f8d221a22e}.
\href{https://github.com/QEEDDD/Verifiable-Mathematical-Reasoning-Corpus}
{Repository and manuscript source}. The necessary flow and covering
arguments are reproduced in this note.

\bibitem{TransferTwo}
\emph{Exact lattice covering transfers: covering minima, partitions,
and arithmetic families}, research manuscript prepared for Keio,
20 November 2026, revised 23-page version. Unpublished preceding note.

\bibitem{TransferThree}
\emph{Lattice covering flows: collinear interior points, higher minima,
and flatness obstructions}, research manuscript prepared for Keio,
20 November 2026, 27-page version. Unpublished preceding note.
\end{thebibliography}
\end{document}
